Parameterized coinduction, or \code{paco} (Hur et al.), is a library for
low-level coinductive reasoning in the Rocq Proof Assistant which greatly
improves on Rocq's awkward \code{cofix} tactic. But proofs using \code{paco}
still involve intricate definitions and a wide array of library tactics that can
make them long and complex. An alternative to \code{paco} is offered by Pous's
\code{rocq-coinduction} library, which mechanizes the chain construction and
tower induction principle of Schäfer and Smolka. Tower induction promises
coinductive proofs that rely on fewer tactics than \code{paco}, and unlike
\code{paco} provides a uniform mechanism for proofs by coinduction and proofs of
soundness of up-to techniques (Pous \& Sangiorgi 2011, Schäfer \& Smolka 2017).

We assessed that promise at scale by porting the coinductive-proof
infrastructure of the Interaction Trees (ITrees) library, a large library that
relies extensively on coinductive reasoning, from \code{paco} to
\code{rocq-coinduction}. This first port was disappointing. Proof size went
down, but tactic surface went up far more. Proof text fell $12\%$, from $400$k to
$352$k characters, while the custom tactic code ITrees had to write in order to
drive the library nearly tripled, a rise of $178\%$. In particular,
\code{rocq-coinduction} required much client-side plumbing around its exposed
tactics in order for them to function, and also contained several usability bugs
that required client-side workarounds in multiple locations. Furthermore,
development occasionally stalled because the underlying tactics in the
tower-based library were written in OCaml and relied on complex under-the-hood
invariants, which made them difficult to debug in practice.

We therefore rewrote the tactic layer of \code{rocq-coinduction} itself from
scratch. The new implementation is Rocq-native, replacing the OCaml plugin with
Gallina and Ltac, and \igap{\tacImplClaim{}} while providing more expressive
tactics and removing the limitations above. Rerunning the ITrees port against it
keeps that proof reduction and removes what it cost. Library tactic invocations
fell $52\%$ ($2705$ to $1307$), the library defines $8$ fewer wrapper tactics on
top of enhanced coinduction compared to \code{paco} ($20$ to $12$), and the
build-time regression the first port introduced fell from $26\%$ to $4\%$. These
results suggest that tower-based coinduction is not only a theoretical
simplification but a practical improvement over \code{paco} for large
coinductive developments.

\note{\textbf{Two numbers changed since the talk, both in your favour, and one
needs a decision.}
\emph{(a)} ``$20$ to $14$'' is now \textbf{$20$ to $12$}, eight fewer rather than
six, because the shared \code{to_mon} driver collapsed two more families. Fixed
above.
\emph{(b)} the \texttt{tacImplClaim} placeholder is a placeholder for the
$65\%$ claim, which depends on a counting convention you have not picked. Slide
83 says ``reduced it by $65\%$'', from ``$-300$ LOC of ocaml, $-200$ LOC of
Rocq''. Counting like for like, comments and blanks removed on both sides, the
tactic implementation goes $363 \to 192$ lines ($-47\%$) or $12\,233 \to 6\,164$
characters ($-50\%$); counting the whole of \code{tactics.v} against the plugin
by raw \code{wc -l} gives $-62\%$. \Cref{tab:library} has the full breakdown.
Pick one and I will wire it through the abstract, the macros and the checker.}
